\chapter{Application1: Climate change and winter wheat yields: the role of freezing days and long-run adaptations}

\timestamp{2026-11-11}{2026-11-11}

\section{这篇文章在说什么}
本论文在估计 climate change impacts on Chinese winter wheat yields的过程中使用了season-specific climate variables和climate penalty mode，得出了两个结论：1、在实证估计中纳入季节性变量是必要的，本研究中Spring freezing days的effects与传统估计完全不同；2、长期适应行为欸也会对climate change effect产生缩小甚至逆转的效果。

\section{数据来源}
\subsection{数据规模：region \times period}
\subsection{yield data}
\subsection{climate data}

\section{temperature thresholds estimate}
Step1: Seasonal division

Step2: Variables construction
figure 1

\begin{card}
    Question1: sum of four variables in given one season is constant，似乎会出现完全共线性问题？ 先插个眼。
\end{card}

Step3: Range loop
table 1

\begin{card}
    "pick the one which has the best fit"

    What's the fit function? 论文正文中并没有，从Appendix或者reference中寻找补充到此处。
\end{card}

\section{The estimate model}
\subsection{The conventional fixed effects panel data model}
equation (1) in paper，有seasonal variables，只是没有添加climate penalty terms

\subsection{The climate penalty model}
A simple climate penalty model: equation (2) in paper

how to interpret coefficients: \beta_1 代表 long-response， (\beta_1, \beta_2)代表 short-response。

\begin{card}
    Trending Historical change，将Appendix中的内容补充到这里。
\end{card}

\section{Empirical results}
Table 2
\subsection{Column 1: Conventional panel approach}
解释系数，并强调这些系数是long-response.

\subsection{COlumn 2/3: Panel approaches with climate penalty terms}
解释系数，
contrast with Column1
Interesting findings and interprations

\subsection{Robustness check}

\subsection{Projecting}
包含 5.4 和5.5 部分。